% --- Chapter: Console Input and Output ---
\section{Console Input and Output}

CRISP provides type-safe console I/O with proper flush semantics. Output functions support both bare syntax (Perl-style) and parenthesized syntax (Python-style).

\subsection{Output Functions}

\begin{tcolorbox}[colback=codebg, colframe=darkpurple, colbacktitle=darkpurple, title=\textbf{\textcolor{codegold}{Output Functions}}, fonttitle=\bfseries, sharp corners]
\begin{lstlisting}
# print — output without newline (bare syntax)
print "Hello, ";
print "World!";
# Output: Hello, World!

# print — output without newline (paren syntax)
print("Hello, ", "World!");
# Output: Hello, World!

# say — output with automatic newline (bare syntax)
say "Hello, CRISP!";
# Output: Hello, CRISP!\n

# say — output with automatic newline (paren syntax)
say("Hello, ", "CRISP!");
# Output: Hello, CRISP!\n

# Multiple arguments
say "Name: ", name, ", Age: ", age;

# warn — output to stderr
warn("Something might be wrong");
# Output (to stderr): Warning: Something might be wrong
\end{lstlisting}
\end{tcolorbox}

\subsection{Input Functions}

Three input functions are available, each with different behavior:

\begin{table}[H]
\centering
\rowcolors{2}{softviolet!30}{white}
\caption{\textcolor{darkpurple}{Input Functions}}
\vspace{0.3em}
\begin{tabular}{@{}lllp{6cm}@{}}
\toprule
\rowcolor{softvioletdark}
\textbf{\textcolor{white}{Function}} & \textbf{\textcolor{white}{Returns}} & \textbf{\textcolor{white}{Flushes?}} & \textbf{\textcolor{white}{Description}} \\
\midrule
\texttt{readline(prompt?)} & Always String & Yes & Reads a line, strips trailing newline \\
\texttt{read(prompt?)} & Int, Float, or String & Yes & Auto-detects numeric types \\
\texttt{input(prompt?)} & Always String & Yes & Python-compatible alias for readline \\
\bottomrule
\end{tabular}
\end{table}

\begin{tcolorbox}[colback=codebg, colframe=darkpurple, colbacktitle=darkpurple, title=\textbf{\textcolor{codegold}{Input Functions}}, fonttitle=\bfseries, sharp corners]
\begin{lstlisting}
# readline — always returns String
let name = readline("What's your name? ");
say "Hello, ", name, "!";

# read — auto-detects Int, Float, or String
let age = read("How old are you? ");
match age {
    n where type(n) == "int" => { say "Age: ", n; }
    n where type(n) == "float" => { say "Weird age: ", n; }
    _ => { say "That's not a number!"; }
}

# input — Python-compatible alias
let color = input("Favorite color? ");
say "You like ", color;

# Empty input (just press Enter) returns empty string
let data = readline("> ");
if data == "" {
    say "No input provided.";
}

# EOF (Ctrl+D) returns empty string
let eof = readline();
if eof == "" {
    say "Goodbye!";
}
\end{lstlisting}
\end{tcolorbox}

\subsection{Flush Semantics}

All input functions flush stdout before blocking on stdin. This guarantees the prompt is visible before the program waits for input:

\begin{tcolorbox}[colback=codebg, colframe=darkpurple, colbacktitle=darkpurple, title=\textbf{\textcolor{codegold}{Flush Guarantee}}, fonttitle=\bfseries, sharp corners]
\begin{lstlisting}
# This works correctly — prompt appears before blocking
let name = readline("Name: ");
# User sees: Name: _ (cursor waits here)

# Without flush, the prompt might be buffered and invisible
# CRISP guarantees flush, so this never happens
\end{lstlisting}
\end{tcolorbox}

\subsection{Type Conversion Functions}

\begin{tcolorbox}[colback=codebg, colframe=darkpurple, colbacktitle=darkpurple, title=\textbf{\textcolor{codegold}{Type Conversion}}, fonttitle=\bfseries, sharp corners]
\begin{lstlisting}
# int() — parse string to integer, returns null on failure
let n = int("42");
say n;        # 42 (Int)

let bad = int("hello");
say bad;      # null

# float() — parse string to float, returns null on failure
let f = float("3.14");
say f;        # 3.14 (Float)

let bad2 = float("world");
say bad2;     # null

# type() — get the type of any value
say type(42);        # "int"
say type("hello");   # "string"
say type([1, 2]);    # "array"
say type(null);      # "null"
say type(true);      # "bool"

# len() — length of collection or string
say len("CRISP");                  # 5
say len([1, 2, 3, 4, 5]);         # 5
say len({ a => 1, b => 2 });      # 2
\end{lstlisting}
\end{tcolorbox}
